\documentclass{article}
\usepackage[T1]{fontenc}
\usepackage{lmodern}
\usepackage{graphicx}
\usepackage{amsmath,amssymb}
\usepackage[a4paper,margin=2cm]{geometry}
\usepackage[absolute,overlay]{textpos}
\setlength{\TPHorizModule}{1mm}
\setlength{\TPVertModule}{1mm}
\usepackage[indonesian]{babel}
\usepackage{hyperref}
\setlength{\emergencystretch}{2em}
% unit: Halaman 5

\begin{document}

% === Halaman 5 (llm) ===

\textbf{Contoh 1:} Misalkan $n = 6$ dan $w_1 = 1$, $w_2 = 5$, $w_3 = 6$, $w_4 = 11$, $w_5 = 14$, dan $w_6 = 20$.

\textbf{Plainteks:} $\color{red}{1110010101100000000011000}$

Plainteks dibagi menjadi blok yang panjangnya 6, kemudian setiap bit di dalam blok dikalikan dengan $w_i$ yang berkoresponden sesuai dengan persamaan (1):

\begin{align*}$
\text{Blok plainteks ke-1} & : 111001 \\ 
\text{Kriptogram} & : (1 \times 1) + (1 \times 5) + (1 \times 6) + (0 \times 11) + (0 \times 14) + (1 \times 20) = 32 \\
\\
\text{Blok plainteks ke-2} & : 010110 \\ 
\text{Kriptogram} & : (1 \times 5) + (1 \times 11) + (1 \times 14) = 30 \\
\\
\text{Blok plainteks ke-3} & : 000000 \\ 
\text{Kriptogram} & : 0 \\
\\
\text{Blok plainteks ke-4} & : 011000 \\ 
\text{Kriptogram} & : (1 \times 5) + (1 \times 6) = 11
\end{align*}

Jadi, cipherteks yang dihasilkan: $\color{red}{32 \ 30 \ 0 \ 11}$

\end{document}
